\begingroup
\fancyhead[L]{\small NACHLASS UND BRIEFWECHSEL}
\clearpage
\SourceLeaf{571}
\pdfbookmark[1]{Nachlass VII. Tafeln für die Störungen der Pallas durch Jupiter}{pallas-nachlass-VII}
\begin{center}
\large NACHLASS UND BRIEFWECHSEL.
\par\bigskip
\Large [VII. TAFELN FÜR DIE STÖRUNGEN\par DER PALLAS DURCH JUPITER.]
\par\medskip
\normalsize [AUGUST 1816---DEZEMBER 1817.]
\par\bigskip
[1.]
\par\medskip
\textsc{Gauss} an \textsc{Encke}, Göttingen, 5. August 1816.
\end{center}
\pdfbookmark[2]{[1] Gauss an Encke, 5. August 1816}{pallas-VII-letter1}
\par\noindent\hspace*{\parindent}.\ .\ .\ .\ .\quad Zugleich habe ich noch eine andere Bitte. Die gesamte Rechnung
aller Störungen der Pallas für eine einzige Opposition dauert immer
einen Monat, weil man die Rechnungen hin- und herschicken muss. Ich selbst
habe (zwar nicht ununterbrochen, aber man hat auch selten Lust dazu)
diesmal fast 14 Tage mit der einmaligen Rechnung zugebracht. Ich habe hin
und her überlegt, wie man die Arbeit durch Tafeln abkürzen könnte,
und gefunden, dass man sie mithilfe einer gewissen Tafel tatsächlich
vielleicht in 3 Stunden erledigen könnte. Die Berechnung der Tafel wird
aber ungeheure Arbeit machen. Vor allem müsste man abschätzen, wie
\textls[100]{viel} Arbeit es wäre, um beurteilen zu können, ob vielleicht 6 oder 8 Rechner
die Aufgabe ausführen könnten, wenn auch erst nach Jahren. Zu diesem Zweck
schicke ich Ihnen hier die Daten für eine Proberechnung. Mir liegt
weniger an den Ergebnissen (wenigstens vorerst) als daran, genau zu wissen,
wie viel Zeit Sie dafür etwa gebraucht haben werden. Diese Rechnung dürfte
etwa \(\tfrac{1}{100}\) der \textls[100]{gesamten} Arbeit ausmachen.
\par\noindent\hspace*{\parindent}Ich möchte die Werte haben
\SourceLeaf{572}
\[
\begin{aligned}
\text{1) von }P={}&+359\overset{\prime\prime}{,}834\\
&+54{,}481\cos x &{}+48\overset{\prime\prime}{,}391\sin x\\
&+14{,}295\cos2x &{}-36{,}951\sin2x\\
&+16{,}161\cos3x &{}-8{,}251\sin3x\\
&+6{,}273\cos4x &{}+2{,}923\sin4x\\
&+0{,}162\cos5x &{}+6{,}354\sin5x\\
&-0{,}674\cos6x &{}+3{,}129\sin6x\\
&+1{,}702\cos7x &{}+0{,}899\sin7x\\
&+0{,}922\cos8x &{}-0{,}065\sin8x\\
&+0{,}460\cos9x &{}-0{,}414\sin9x\\
&+0{,}181\cos10x &{}-0{,}326\sin10x\\
&+0{,}063\cos11x &{}-0{,}198\sin11x\\
&+0{,}023\cos12x &{}-0{,}101\sin12x
\end{aligned}
\]
\[
\begin{aligned}
\text{2) von }Q={}&-544\overset{\prime\prime}{,}924\\
&-98{,}084\cos x &{}-6\overset{\prime\prime}{,}178\sin x\\
&+27{,}589\cos2x &{}-10.928\sin2x\\
&+8{,}342\cos3x &{}+4{,}711\sin3x\\
&-0{,}875\cos4x &{}+1{,}294\sin4x\\
&-4{,}480\cos5x &{}-1{,}766\sin5x\\
&-1{,}813\cos6x &{}-1{,}284\sin6x\\
&-0{,}179\cos7x &{}+1{,}643\sin7x\\
&+0{,}434\cos8x &{}+1{,}020\sin8x\\
&+0{,}414\cos9x &{}+0{,}460\sin9x\\
&+0{,}326\cos10x &{}+0{,}181\sin10x\\
&+0{,}198\cos11x &{}+0{,}063\sin11x\\
&+0{,}101\cos12x &{}+0{,}023\sin12x
\end{aligned}
\]
\noindent\hspace*{\parindent}3) von \(\log N\) und \(\psi\), sodass
\[
P=N\cos\psi,\qquad Q=N\sin\psi.
\]
\noindent\(\psi\) wird am Ende auf \textls[100]{Bogensekunden}, \(\log N\) auf 5 Dezimalstellen verlangt. Wünschenswert
wäre diese Tafel zwar für \(x=0\), \(x=1^0\), \(x=2^0\), \(x=3^0\) usw. bis
\(x=360^0\), aber das wäre eine viel zu gewaltige Arbeit. Man
\SourceLeaf{573}
muss sich also stärker beschränken. Ich denke, dass man sich \textls[100]{allenfalls} damit begnügen kann,
die Werte von \(x\) von \(15^0\) zu \(15^0\) wachsen zu lassen, sodass
24 solche Rechnungen auszuführen sind. So möchte ich daher die Proberechnung ausgeführt haben.
Dabei kann man übrigens einige Vorteile in einzelnen Fällen nutzen,
die Sie leicht selbst finden werden und die von
\[
\left.\begin{aligned}
\sin15^0&=\BareRoot{\tfrac{3}{8}}-\BareRoot{\tfrac{1}{8}}\\
\sin45^0&=2\,\mathord{.}\,\BareRoot{\tfrac{1}{8}}
&\qquad\sin60^0&=\BareRoot{\tfrac{3}{4}}\\
\sin75^0&=\BareRoot{\tfrac{3}{8}}+\BareRoot{\tfrac{1}{8}}
\end{aligned}\right\}\ \text{abhängen.}
\]
\noindent\hspace*{\parindent}.\ .\ .\ .\ .
\clearpage
\endgroup
